\documentclass[../main.tex]{subfiles}
\begin{document}

\section{Summary}

Luận văn xuất phát từ một hạn chế cụ thể của các kiến trúc agent hiện có cho ISD: chúng chủ yếu tối ưu cho chất lượng artifact tương thích ADDIE tại lần sinh đầu tiên, chưa đặt mức độ alignment giữa các thành phần thiết kế làm trạng thái trung tâm có thể truy vết, kiểm tra, sửa chữa hay tinh chỉnh hiệu quả cho nhu cầu cải tiến liên tục. Để giải quyết hạn chế đó, luận văn đề xuất AlignmentGraph-ISD, một alignment-graph harness biểu diễn objective, assessment, activity, material, implementation plan và evaluation cùng các quan hệ \texttt{measures}/\texttt{supports}/\texttt{evaluates} giữa chúng dưới dạng một graph tường minh, được vận hành bởi một dynamic ISD task graph, một tập agent chuyên biệt theo vai trò, cơ chế schema validation/pedagogical verification/targeted repair kèm cascade invalidation, và một bộ nhớ bốn lớp.

Khi được tích hợp và đánh giá như một agent mới trong ISD-Agent-Bench, AlignmentGraph-ISD cho kết quả nhất quán trên hai lượt benchmark dùng hai agent model khác nhau: đứng thứ hai trên subset sơ bộ \texttt{test\_30} (85.67 điểm trung bình, sau React-ADDIE) và dẫn đầu trên subset mở rộng \texttt{test\_90} (84.07 điểm trung bình, dẫn đầu ADDIE score và cả ba mức difficulty), với tỷ lệ thành công 100\% trên cả hai lượt chạy. Khoảng cách còn tồn tại rõ nhất là trajectory score, nơi AlignmentGraph-ISD vẫn thấp hơn React-ADDIE ở cả hai lượt, cho thấy cơ chế logging và dynamic orchestration của harness còn có thể cải thiện thêm. Bên cạnh kết quả định lượng, luận văn trình bày thiết kế một case study trong miền IT/Software Engineering, minh hoạ cách core harness được mở rộng thành một ứng dụng có quản lý phiên nhiều lượt, truy xuất tri thức theo domain và giao diện teacher-facing; phần lớn thiết kế này, đặc biệt là UI và cơ chế export, tại thời điểm hoàn thành luận văn vẫn ở trạng thái dự kiến, chưa được hiện thực đầy đủ.

Nhìn chung, luận văn đã trả lời được RQ1 với bằng chứng nhất quán trên hai cỡ mẫu và hai model khác nhau, cho thấy AlignmentGraph-ISD cải thiện chất lượng artifact so với các agent baseline và các ISD agent hiện có. RQ2 mới được trả lời ở mức cơ chế và giao thức: chuỗi verifier--repair--cascade làm nền cho khả năng tinh chỉnh đã vận hành ổn định bên trong từng lượt chạy benchmark và được ghi trong trajectory, giao thức đo với bốn process metric (cascade scope, edit locality, cost per tuning round, alignment preservation) đã được định nghĩa ở Chương 4, nhưng các phiên tinh chỉnh của case study chưa được thực thi nên chưa có số liệu nhiều phiên. Tương tự, các ablation (\texttt{no\_verifier}, \texttt{no\_repair}, \texttt{single\_agent}, \texttt{linear\_addie\_only}) — phân tích bổ trợ về đóng góp của từng thành phần harness — đã được thiết kế đầy đủ nhưng chưa có số liệu thực nghiệm tương ứng.

\section{Suggestion for Future Works}

Hướng phát triển gần nhất là hoàn thành các phần còn dở của luận văn theo đúng thiết kế đã trình bày ở Chương 4 và Chương 5: thực thi các phiên tinh chỉnh của case study để đo bốn process metric và trả lời trọn vẹn RQ2; chạy đầy đủ bốn ablation đã thiết kế để định lượng đóng góp của từng thành phần harness; mở rộng đánh giá định lượng từ \texttt{test\_90} lên full test split 1202 scenario để tăng độ tin cậy thống kê của so sánh; và hoàn thiện case study IT/Software Engineering với đầy đủ giao diện teacher-facing, cơ chế export Markdown/PDF, và một đợt expert review thật thay vì chỉ dừng ở thiết kế tiêu chí đánh giá.

Về lâu dài, khoảng cách trajectory score so với React-ADDIE gợi ý hướng cải thiện cơ chế logging và dynamic orchestration của \texttt{MetaAgent}, để trace phản ánh đầy đủ hơn quá trình ra quyết định thay vì chỉ ghi lại kết quả cuối của mỗi bước. Một hướng mở rộng khác là kiểm tra khả năng khái quát hoá của AlignmentGraph-ISD sang các domain ngoài IT/Software Engineering, thông qua việc thay corpus, template và verifier policy mà không cần thay đổi core harness, nhằm kiểm chứng mức độ đúng của giả định rằng alignment graph là một cấu trúc trung lập với domain.

\end{document}